The variational theory of the one-phase Bernoulli problem begins with
Alt and Caffarelli's existence and regularity theory
\cite{AltCaffarelli1981}.  Weiss's monotonicity formula and dimension
reduction \cite{Weiss1999} connect the singularities of general minimizers
to homogeneous global minimizers; see also
\cite[Theorem~1.4 and Chapters~9--10]{Velichkov2023}.
Thus the first dimension admitting a nonflat minimizing cone determines
both the threshold for singular free boundaries and the codimension
bound furnished by this theory.  Our main result answers positively
the critical-dimension question whether $d_*=7$.

A second strand of the theory passes from weak free boundaries to
classical ones. Caffarelli's Harnack-inequality approach established
regularity for Lipschitz free boundaries \cite{Caffarelli1987}.
De Silva developed a different improvement-of-flatness argument allowing
a nonzero right-hand side \cite[Theorems~1.1--1.2]{DeSilva2011}; this is the approach
presented in the regularity chapters of \cite{Velichkov2023}.
Higher regularity has classical antecedents in Kinderlehrer and Nirenberg
\cite[Theorem~2]{KinderlehrerNirenberg1977}. The version used here, for
Lipschitz viscosity solutions, is due to De Silva, Ferrari and Salsa
\cite[Theorem~1.1 and Corollary~1.2]{DeSilvaFerrariSalsa2019}.
These results make a classification of
homogeneous minimizing blowups applicable to general local minimizers.

As recalled by Jerison and Savin \cite[Introduction]{JerisonSavin2015},
the dimension-three rigidity result for stable homogeneous solutions
with smooth cross-section is due to Caffarelli, Jerison and Kenig
\cite{CaffarelliJerisonKenig2004}.
Jerison and Savin proved that stable homogeneous solutions with smooth
cross-section are flat in dimensions at most four
\cite[Theorem~1.1]{JerisonSavin2015}; their regularity consequence for
minimizers \cite[Corollary~1.2]{JerisonSavin2015} gives $d_*\geq5$.
Their argument constructs functions of the Hessian satisfying
inequalities for the linearized free-boundary problem.  In the other
direction, De Silva and Jerison constructed a nonflat homogeneous
global minimizer in dimension seven
\cite[Theorem~1.1]{DeSilvaJerison2009}, proving $d_*\leq7$.
The minimizing property in that construction is with respect to all
Sobolev competitors with the prescribed boundary trace. Firester,
Tsiamis and Wang also construct further symmetric minimizing
one-phase cones, including a second example in dimension seven
\cite[Theorem~1.4]{FiresterTsiamisWang2026Capillary}.

Recent work further describes cones with symmetry.  Firester, Tsiamis
and Wang \cite[Theorem~1.1]{FiresterTsiamisWang2026} prove that their
cohomogeneity-one cones with $O(n-k)\times O(k)$ symmetry,
$1\leq k\leq n-2$, are unstable for $n\leq6$ and strictly stable for
$n\geq7$.  The exclusion proved here applies to globally
minimizing cones without a symmetry assumption.

Cone classification also differs from rigidity of arbitrary entire
stable solutions. Fern\'andez-Real and Serra prove that every global
classical stable solution in dimension four is one-dimensional
\cite[Theorem~1.1]{FernandezRealSerra2026}. Their result removes homogeneity in
that dimension.

For an admissible increasing family of nonnegative boundary data,
Fern\'andez-Real and Yu prove regularity through dimension six for
almost every parameter
\cite[Assumption~1 and Theorem~1.3]{FernandezRealYu2025}.
This is a generic-data result: the increase is required at a fixed linear
rate on the positive set of the smaller boundary datum.

\paragraph{The obstacle and the proof strategy.}
The stability inequality compares the Dirichlet energy of a test function
with a curvature term on the free boundary.  The use of the harmonic Hessian has a geometric precedent in Simons's
second-fundamental-form argument for minimal cones \cite[Section~6.1]{Simons1968},
a connection developed explicitly by Jerison and Savin. A natural
Bernoulli test is a power of the Hessian norm. Jerison and Savin explain why
this unweighted choice does not give a uniform instability criterion in
dimensions four through six: the interior and boundary estimates impose
different restrictions on the power, and their sharp configurations need
not agree
\cite[Section~3.2]{JerisonSavin2015}.  Their four-dimensional argument
uses a more refined function of the Hessian eigenvalues
\cite[Section~4]{JerisonSavin2015}.

Our proof uses the normalized shape of the Hessian to construct both
a scalar test and a separate vector field whose divergence controls
the test's energy.  The two constructions are chosen together, so
the interior and boundary estimates retain the same tensor data.
The proof has three stages.

First, suppose the first nonflat minimizing cone occurs in dimension
five or six.  Blowups away from its vertex reduce the dimension and
are therefore flat.  Regularity theory makes each component of its
spherical positive set smooth, and at least one component has
nonzero Hessian.  A five-dimensional cone can be made independent
of one extra coordinate without losing minimality.  Its spherical
link then has precisely two exceptional tips.  Thus both possible
dimensions lead to a domain $\Omega\subset S^5$ with a smooth
regular boundary and, possibly, these two tips
(Section~\ref{sec:reductions}).

Second, let $g$ be the restriction of the resulting cone to the sphere,
and set $T=\nabla^2g+g\Id$ and $w=|T|$.  Where $w>0$, we construct
a positive scalar $h$ from $w$, the normalized tensor $T/w$, $g$,
and $\nabla g$, and
a vector field $Z$ from the same data and $\nabla T/w$.  The target
inequalities are
\[
  \diver Z\ge |\nabla h|^2+4h^2+\epsilon_*w^{2a},
  \qquad N\cdot Z+Hh^2\ge0,
\]
where $a\in(1/2,1)$ is fixed, $\epsilon_*>0$, $N$ points into $\Omega$, and
$H=-T(N,N)$.  Spherical stability supplies the reverse integrated
comparison between the test's energy and its boundary curvature term.
The extra positive mass $\epsilon_*w^{2a}$ will give the contradiction.
The geometry and stability inequality are developed in
Section~\ref{sec:geometry}. We then state the universal tensor inequality
and derive the two geometric comparisons; the full exact verification
follows in Section~\ref{sec:certificate-proof}.

Two cancellations make this construction possible.  A Laplacian
contraction and an antisymmetric curvature contraction eliminate
fourth derivatives from $\diver Z$.  At the boundary, reflection
in the normal direction removes the derivative components not
determined by the boundary conditions.  The remaining inequalities
involve only finitely many normalized tensors.  We prove them by
weighted sums of tensor squares and exact rational remainder bounds,
with all coefficients and contraction conventions supplied in
Appendix~\ref{app:tables}.

Finally, neither the normalized Hessian nor its derivatives are defined
at $w=0$, and suspension introduces singular tips.  These sets cannot
be discarded without an estimate.  A regularized Bochner identity
gives an integrable weight controlling the derivatives near every
Hessian zero.  The explicit suspension formulas give an integrable
power of the distance parameter at both tips.  Smooth product
cutoffs therefore remove both obstructions and make all error terms
vanish.  Section~\ref{sec:cutoffs} completes the stability
contradiction.
